% Generated from locale/tr/content/lambda-calculus/introduction/fixed-point-combinator.tex; do not hand-edit.


\olfileid[tr]{lam}{int}{fix}
\olsection{Sabit Nokta Birleştiricileri}

Bir $g$ lambda teriminiz olduğunu ve $k$'nin $gk$ ile $\beta$-eşdeğer olması
özelliğini taşıyan başka bir $k$ terimi istediğinizi varsayın. Gösterim
uzlaşımlarımızı kullanarak
\[
\fn{diag}(x) = xx
\]
ve
\[
l(x) = g(\fn{diag}(x))
\]
terimlerini tanımlayın; başka bir deyişle $l$,
$\lambd[x][g(xx)]$ terimidir. $k$, $ll$ terimi olsun. O zaman
\begin{align*}
k & = (\lambd[x][g(xx)])(\lambd[x][g(xx)]) \\
& \red  g((\lambd[x][g(xx)])(\lambd[x][g(xx)])) \\
& = gk.
\end{align*}
Eğer
\[
Y = \lambd[g][((\lambd[x][g(xx)])(\lambd[x][g(xx)]))]
\]
alınırsa $Yg$ ile $g(Yg)$ ortak bir terime indirgenir; dolayısıyla
$Yg\equiv_\beta g(Yg)$. Buna ``Curry birleştiricisi'' denir. Bunun yerine
\[
Y = (\lambd[xg][g(xxg)])(\lambd[xg][g(xxg)])
\]
alınırsa, daha güçlü olarak $Yg$ gerçekten $g(Yg)$'ye indirgenir. $Y$'nin
bu ikinci sürümüne ``Turing birleştiricisi'' denir.

